\documentclass[12pt]{article}
\usepackage[utf8]{inputenc}
\usepackage[spanish, es-noshorthands]{babel}
\decimalpoint
\usepackage{amssymb,amsmath,amsthm,amsfonts,mathtools}
\usepackage[breakable]{tcolorbox}
\usepackage{microtype}

% ---- COLORES ----
\usepackage{xcolor}
\definecolor{miAzul}{RGB}{13, 33, 161}
\definecolor{miRojo}{RGB}{186, 24, 27}
\definecolor{miVerde}{RGB}{65, 123, 90}

% ---- LINKS ----
\usepackage{hyperref}
\hypersetup{colorlinks=true, urlcolor=miAzul, citecolor=miVerde, linkcolor=miRojo}

% ---- IMÁGENES ----
\usepackage{graphicx}
\graphicspath{ {imgs/} }

% ---- FORMATO DE TÍTULOS ----
\usepackage{titlesec}
\titlespacing{\section}{0pt}{1em plus 4pt minus 2pt}{0pt plus 2pt minus 2pt}
\titlespacing{\subsection}{0pt}{1em plus 4pt minus 2pt}{0pt plus 2pt minus 2pt}
\titlespacing{\subsubsection}{0pt}{1em plus 4pt minus 2pt}{0pt plus 2pt minus 2pt}

% ---- FORMATO DE LISTAS ----
\usepackage[inline]{enumitem}
\setlist[itemize]{topsep=-0.5em} % es media m para compensar
\setlist[enumerate]{topsep=-0.5em}

% ---- GRÁFICAS ----
\usepackage{tikz}
\usepackage{scalerel}
\usepackage{pict2e}
\usepackage{tkz-euclide}
\usetikzlibrary{calc}
\usetikzlibrary{patterns,arrows.meta}                   % Mejores flechas
\usetikzlibrary{shadows}                                % Sombras
\usetikzlibrary{external}                               % ¿?
\usetikzlibrary{decorations.pathreplacing,calligraphy}  % Para poder graficar el '{'
% Plots
\usepackage{pgfplots}
\pgfplotsset{compat=newest}                             % Ancho de las gráficas: width=6.5cm,
\usepgfplotslibrary{statistics}
\usepgfplotslibrary{fillbetween}

% ---- BIBLIOGRAFÍA ----
\usepackage[backend=biber]{biblatex}
\addbibresource{fuentes.bib}

% ---- PSEUDOCÓDIGO ----
\usepackage{algorithm}
\usepackage{algpseudocode}
\usepackage{amsmath}

% ---- FORMATO DE LAS PÁGINAS ----
\usepackage[a4paper, width=165mm, top=20mm, bottom = 20mm]{geometry}
\usepackage{fancyhdr}
\pagestyle{fancy}
\fancyhf{}
\lhead{\footnotesize Proyecto final}
\chead{\footnotesize \textbf{Compresión de Imágenes con SVD}}
\rhead{\footnotesize \today}
%\lfoot{\footnotesize Equipo 2}
\cfoot{\footnotesize \textsc{Análisis Aplicado I}}
\rfoot{\footnotesize \thepage}
\renewcommand{\headrulewidth}{0.4pt}
\renewcommand{\footrulewidth}{0.4pt}
\setlength{\headheight}{15pt}

% Espacio entre párrafos y quitar sangría
\setlength{\parskip}{1em}
\setlength{\parindent}{0cm}

\begin{document}
\begin{titlepage}
	\centering
	\includegraphics[width=0.5\textwidth]{logo-ITAM.pdf}\par\vspace{1cm}
	{\scshape\large Instituto Tecnológico Autónomo de México \par}
	\vspace{2.5cm}
	
	{\scshape\Large Análisis Aplicado I\par}
	\vspace{1.5cm}
	{\LARGE\bfseries Proyecto final\par}
    \vspace{0.5cm}
	{\large\bfseries Compresión de Imágenes con SVD}
	\vspace{3cm}
	
	{\large
	\begin{tabular}{||l c||} 
         \hline
         \multicolumn{1}{||c}{\textbf{Integrantes del equipo}} & \textbf{Clave única} \\
         \hline\hline
         Darío Sotelo & \texttt{XXXXXX} \\ 
         \hline
         María Vara & \texttt{XXXXXX}\\
         \hline
         José Alberto Márquez Luján & \texttt{187917} \\
         \hline                  
    \end{tabular}
    }
	\vfill
	Profesora:\par
	Dra.~María Victoria Chávez Hernández

	\vfill
	{\today\par}
\end{titlepage}

\section{Introducción}
Así se hace una cita: \cite{Bolotin2017}.


% Imprime las referencias
\printbibliography

\end{document}